\subsection{Factors of a Semisimple Ring}

In this subsection, we consider a semisimple ring A.

We denote the set of classes of simple left A-modules by \(\mathcal{S}\) (VIII, p.~51); it is finite (VIII, p.~136, Proposition 1). For every \(\lambda \in \mathcal{S}\) , we choose a simple A-module \(S_{\lambda}\) of class \(\lambda\) . We denote its annihilator by \(\mathfrak{b}_{\lambda}\) and the opposite field of its commutant \(\mathrm{End}_{\mathrm{A}}(S_{\lambda})\) by \(D_{\lambda}\) .

PROPOSITION 8. --- a) For every \(\lambda \in \mathcal{S}\) , \(S_{\lambda}\) is a finite-dimensional right vector space over the field \(D_{\lambda}\) . When passing to the quotient, the mapping \(a \mapsto a_{S_{\lambda}}\) defines a ring isomorphism from \(A / \mathfrak{b}_{\lambda}\) to \(\mathrm{End}_{D_{\lambda}}(S_{\lambda})\) .

\begin{enumerate}
\def\labelenumi{\alph{enumi})}
\setcounter{enumi}{1}
\tightlist
\item
  For every \(\lambda \in \mathcal{S}\) , the ring \(\mathrm{A} / \mathfrak{b}_{\lambda}\) is simple, and the canonical homomorphism \(\psi\) from \(\mathrm{A}\) to \(\prod_{\lambda \in \mathcal{S}}\mathrm{A} / \mathfrak{b}_{\lambda}\) is a ring isomorphism.
\end{enumerate}

The ring \(A/b_{\lambda}\) is isomorphic to \(A_{S_{\lambda}}\) . By Lemma 1 of VIII, p.~139, this ring is simple. The given mapping from \(A/b_{\lambda}\) to \(\operatorname{End}_{D_{\lambda}}(S_{\lambda})\) can be identified with the mapping from \(A_{S_{\lambda}}\) to \(A_{S_{\lambda}}^{\prime\prime}\) . By Proposition 5 of VIII, p.~139, it is an isomorphism.

The A-module \(A_{s}\) is semisimple, faithful, and balanced. The homomorphism \(\psi\) can be identified with the morphism from the bicommutant of \(A_{s}\) to \(\prod_{\lambda\in\mathcal{S}}\mathrm{End}_{\mathrm{D}_{\lambda}}(\mathrm{S}_{\lambda})\) , which is an isomorphism (VIII, p.~86, Proposition 7, c)).

The simple ring \(A/b_{\lambda}\) is called the simple factor of A of type \(\lambda\) .

Example. --- Let K be an algebraically closed commutative field, and let A be a semisimple algebra of finite degree over K. Let \((\mathrm{V}_{i})_{i\in\mathrm{I}}\) be a family of simple A-modules such that every simple A-module is isomorphic to a unique \(V_{i}\) . Then I is a finite set (VIII, p.~136, Proposition 1), the vector spaces \(V_{i}\) are finite-dimensional over the field K, the commutant of \(V_{i}\) is equal to \(K\cdot1_{V_{i}}\) (VIII, p.~47, Theorem 1), and the mapping \(a\mapsto(a_{V_{i}})_{i\in\mathrm{I}}\) is an algebra isomorphism from A to \(\prod_{i\in\mathrm{I}}\mathrm{End}_{\mathrm{K}}(\mathrm{V}_{i})\) (Proposition 8).

We have defined (VIII, p.~50) a minimal two-sided ideal as a minimal element of the set of nonzero two-sided ideals, ordered by inclusion. In other words, a minimal two-sided ideal \(\mathfrak{a}\) of A is a simple (A, A)-sub-bimodule of \(s\mathrm{A}_d\) . Likewise, we define a maximal two-sided ideal of A as a maximal element of the set of proper two-sided ideals of A. A maximal two-sided ideal \(\mathfrak{a}\) of A is simply a maximal (A, A)-sub-bimodule of \(s\mathrm{A}_d\) (VIII, p.~48, Definition 2). If the ring A is simple, then the ideal 0 is a maximal two-sided ideal of A and A is a minimal two-sided ideal.

For any \(\lambda \in \mathcal{S}\) , we denote the isotypical component of type \(\lambda\) of the A-module \(A_s\) by \(\mathfrak{a}_{\lambda}\) . For any subset \(\Lambda\) of \(\mathcal{S}\) , we set \(\mathfrak{a}_{\Lambda} = \sum_{\lambda \in \Lambda} \mathfrak{a}_{\lambda}\) .

PROPOSITION 9. --- a) Order the set \(\mathfrak{P}(\mathcal{S})\) of subsets of \(\mathcal{S}\) and the set \(\mathcal{B}_{\mathrm{A}}\) of two-sided ideals of A by inclusion. The mapping \(\Lambda \mapsto \mathfrak{a}_{\Lambda}\) is an isomorphism of ordered sets from \(\mathfrak{P}(\mathcal{S})\) to \(\mathcal{B}_{\mathrm{A}}\) .

\begin{enumerate}
\def\labelenumi{\alph{enumi})}
\setcounter{enumi}{1}
\item
  The minimal two-sided ideals of A are the ideals \(\mathfrak{a}_{\lambda}\) .
\item
  We have \(\mathfrak{b}_{\lambda} = \mathfrak{a}_{\mathcal{S} - \{\lambda\}}\) for every \(\lambda \in \mathcal{S}\) , and the ideals \(\mathfrak{b}_{\lambda}\) are the maximal two-sided ideals of A.
\item
  For every \(\lambda \in \mathcal{S}\) , the canonical mapping from A to \(\mathrm{A} / \mathfrak{b}_{\lambda}\) induces an isomorphism of A-modules from \(\mathfrak{a}_{\lambda}\) to \(\mathrm{A} / \mathfrak{b}_{\lambda}\) .
\end{enumerate}

Assertion a) follows from Proposition 8, d) of VIII, p.~87 applied to the A-module \(\mathbf{A}_s\) . It follows that the minimal two-sided ideals of A are the \(\mathfrak{a}_{\lambda}\) and that the maximal two-sided ideals are the ideals \(\mathfrak{c}_{\lambda} = \mathfrak{a}_{\mathcal{S} - \lambda}\) (for \(\lambda \in \mathcal{S}\) ).

It remains to show the equality of \(b_{\lambda}\) and \(c_{\lambda}\) for every \(\lambda \in S\) . Let \(\lambda\) and \(\mu\) be distinct in S. The A-submodule \(a_{\mu}S_{\lambda}\) of \(S_{\lambda}\) is the union of the images of the linear mappings \(a \mapsto ax\) from \(a_{\mu}\) to \(S_{\lambda}\) for \(x \in S_{\lambda}\) . Consequently, it is zero, and we have \(a_{\mu} \subset b_{\lambda}\) . We therefore have \(c_{\lambda} \subset b_{\lambda}\) , and finally \(c_{\lambda} = b_{\lambda}\) because \(c_{\lambda}\) is a maximal two-sided ideal of A and \(b_{\lambda}\) is distinct from A.

COROLLARY. --- Let \((\mathrm{A}_{i})_{i\in\mathrm{I}}\) be a finite family of simple rings and f an isomorphism from A to \(\prod_{i\in I}A_{i}\) . For every \(i\in I\) , there exists a unique element \(\varphi(i)\) of S such that the kernel of \(pr_{i}\circ f\) is \(\mathfrak{b}_{\varphi(i)}\) . The mapping \(\varphi\) is a bijection from I to S; the mapping \(pr_{i}\circ f\) induces an isomorphism \(f_{i}\) from \(\mathrm{A}/\mathfrak{b}_{\varphi(i)}\) to \(A_{i}\) for every \(i\in I\) .

Thus, f is the composition of the canonical isomorphism from A to \(\prod_{\lambda\in\mathscr{S}}A/\mathfrak{b}_{\lambda}\) and the isomorphism from \(\prod_{\lambda\in\mathscr{S}}A/\mathfrak{b}_{\lambda}\) to \(\prod_{i\in I}A_{i}\) deduced from the \(f_{i}\) (``uniqueness of the decomposition of a semisimple ring into a product of simple rings'').

Let us prove the corollary. Let \(i \in I\) ; denote the kernel of \(pr_{i} \circ f\) by \(b_{i}^{\prime}\) . Since the simple ring \(A_{i}\) is isomorphic to \(A/b_{i}^{\prime}\) , the two-sided ideal \(b_{i}^{\prime}\) of A is maximal. By Proposition 8, c), there consequently exists a unique element \(\varphi(i)\) of S such that we have \(b_{i}^{\prime} = b_{\varphi(i)}\) . When passing to the quotient, \(pr_{i} \circ f\) defines an isomorphism \(f_{i}\) from \(A/b_{\varphi(i)}\) to \(A_{i}\) . Moreover, we have \(b_{i}^{\prime} + b_{j}^{\prime} = A\) if \(i \neq j\) and \(\cap_{i \in I} b_{i}^{\prime} = 0\) (cf.~I, §8, No.~11, p.~110, Proposition 10). It follows from this and Proposition 8 that \(\varphi\) is a bijection from I to S.

PROPOSITION 10. --- Denote the center of A by Z. For \(\lambda \in \mathcal{S}\) , let \(Z_{\lambda}\) be the center of the field \(D_{\lambda}\) .

\begin{enumerate}
\def\labelenumi{\alph{enumi})}
\item
  The mapping \(z \mapsto (z_{\mathrm{S}_{\lambda}})_{\lambda \in \mathcal{S}}\) is an isomorphism from the ring \(\mathbf{Z}\) to the product \(\prod_{\lambda \in \mathcal{S}}\mathrm{Z}_{\lambda}\) .
\item
  Order the set \(I_{Z}\) of ideals of Z and the set \(B_{A}\) of two-sided ideals of A by inclusion. The mapping \(a \mapsto aA\) is an isomorphism of ordered sets from \(I_{Z}\) to \(B_{A}\) . The inverse isomorphism sends a two-sided ideal b of A to the ideal \(b \cap Z\) of Z.
\end{enumerate}

This proposition follows from Proposition 8 of VIII, p.~87 applied to the A-module \(A_{s}\) , whose bicommutant is A.

COROLLARY. --- Let B be a ring. The following properties are equivalent:

\begin{enumerate}
\def\labelenumi{(\roman{enumi})}
\item
  The ring B is simple.
\item
  The ring B is semisimple, and its center is a field.
\item
  The ring B is semisimple, and there exists only one class of B-simple modules.
\end{enumerate}

PROPOSITION 11. --- Let \(\lambda \in \mathcal{S}\) . The isotypical component \(\mathfrak{a}_{\lambda}\) of A is both the isotypical component of \(\mathrm{A}_s\) of type \(\mathrm{S}_{\lambda}\) and the isotypical component of \(\mathrm{A}_d\) of type \(\mathrm{S}_{\lambda}^{*}\) . Moreover, we have

\[
[ \mathrm{A} _ {s}: \mathrm{S} _ {\lambda} ] = [ \mathrm{A} _ {d}: \mathrm{S} _ {\lambda} ^ {*} ] = \dim_ {\mathrm{D} _ {\lambda}} \mathrm{S} _ {\lambda}\tag{1}
\]

and

\[
\operatorname{long} (\mathrm{A}) = \operatorname{long} \left(\mathrm{A} ^ {\mathrm{o}}\right) = \sum_ {\lambda \in \mathscr {S}} \dim_ {\mathrm{D} _ {\lambda}} \mathrm{S} _ {\lambda}.\tag{2}
\]

The first assertion is the specific case \(M = A_{s}\) of Proposition 6, c) of VIII, p.~139. The equality \([A_{s} : S_{\lambda}] = [A_{d} : S_{\lambda}^{*}]\) follows from Proposition 7 of VIII, p.~140 because the dual of the left A-module \(A_{s}\) is isomorphic to the right A-module \(A_{d}\) . By Propositions 8, a) and 9, c), the mapping \(a \mapsto a_{S_{\lambda}}\) defines an isomorphism of left A-modules from \(a_{\lambda}\) to \(\operatorname{End}_{D_{\lambda}}(S_{\lambda})\) . Since \([A_{s} : S_{\lambda}]\) is, by definition, the length of the left A-module \(a_{\lambda}\) , the relation \([A_{s} : S_{\lambda}] = \dim_{D_{\lambda}} S_{\lambda}\) follows from Lemma 2 of VIII, p.~121. Finally, relation (2) is obtained from (1) by taking the sum over \(\lambda\) .

Scholium. --- Let A be a semisimple ring and Z its center. There exist canonical bijections between the following sets:

\begin{enumerate}
\def\labelenumi{\alph{enumi})}
\item
  The set \(\mathcal{S}(\mathrm{A})\) of classes of simple left A-modules
\item
  The set \(\mathcal{S}(\mathrm{A}^{\circ})\) of classes of simple right A-modules
\item
  The set of minimal two-sided ideals of A
\item
  The set of maximal two-sided ideals of A
\item
  The set \(\mathcal{S}(\mathbf{Z})\) of classes of simple Z-modules
\item
  The set of minimal ideals of Z
\item
  The set of maximal ideals of \(\mathbf{Z}\) .
\end{enumerate}

Thus, to every element \(\lambda\) of \(\mathcal{S}(A)\) , there correspond the class \(\lambda^{*}\) of the simple right A-module \(S_{\lambda}^{*}\) , dual of \(S_{\lambda}\) , the minimal two-sided ideal \(a_{\lambda}\) of A (isotypical component of \(A_{s}\) of type \(\lambda\) ), the maximal two-sided ideal \(b_{\lambda}\) of A (annihilator of the simple module \(S_{\lambda}\) ), the class of the simple Z-module \(Z \cap a_{\lambda}\) , the minimal ideal \(Z \cap a_{\lambda}\) of Z, and the maximal ideal \(Z \cap b_{\lambda}\) of Z.

PROPOSITION 12. --- Let M be a module over the semisimple ring A and \(\mathcal{S}_{\mathrm{M}} \subset \mathcal{S}\) the support of M. Then the annihilator Ann(M) of M is the two-sided ideal \(\sum_{\lambda \in \mathcal{S} - \mathcal{S}_{\mathrm{M}}} \mathfrak{a}_{\lambda}\) , and the trace ideal \(\tau(\mathrm{M})\) of M is the two-sided ideal \(\sum_{\lambda \in \mathcal{S}_{\mathrm{M}}} \mathfrak{a}_{\lambda}\) . In particular, A is the direct sum of Ann(M) and \(\tau(\mathrm{M})\) .

By definition (VIII, p.~84), \(\mathcal{S}_{\mathrm{M}}\) consists of the classes of simple submodules of M. Since the module M is semisimple, the annihilator of M is the intersection of the annihilators \(\mathfrak{b}_{\lambda}\) of the modules of class \(\lambda\) , for \(\lambda\) running through \(\mathcal{S}_{\mathrm{M}}\) . Now, we have \(\mathfrak{b}_{\lambda} = \sum_{\mu \neq \lambda} \mathfrak{a}_{\mu}\) for every \(\lambda \in \mathcal{S}\) (Proposition 9). Since A is the direct sum of the family \((\mathfrak{a}_{\lambda})_{\lambda \in \mathcal{S}}\) , the annihilator of M is indeed equal to \(\sum_{\lambda \in \mathcal{S} - \mathcal{S}_{\mathrm{M}}} \mathfrak{a}_{\lambda}\) .

By definition (VIII, p.~80), the trace ideal \(\tau(\mathrm{M})\) is the A-submodule of \(\mathrm{A}_s\) generated by the images of the A-linear mappings from M to \(\mathrm{A}_s\) . Because M is semisimple, it amounts to the same to say that \(\tau(\mathrm{M})\) is generated by the simple submodules of \(\mathrm{A}_s\) with class belonging to \(\mathcal{S}_{\mathrm{M}}\) . We therefore have \(\tau(\mathrm{M}) = \sum_{\lambda \in \mathcal{S}_{\mathrm{M}}} \mathfrak{a}_{\lambda}\) .

COROLLARY. --- Let M be a module over the semisimple ring A. The following properties are equivalent:

\begin{enumerate}
\def\labelenumi{(\roman{enumi})}
\item
  The A-module M is faithful.
\item
  The support of M is equal to S.
\item
  The A-module M is generating.
\end{enumerate}

Indeed, saying that M is faithful means that its annihilator is reduced to 0, and M is generating if and only if its trace is equal to A (VIII, p.~80, Theorem 1).

